% Chapter 2: Literature Review (the paper plan's "Related Work").
% Source: /mnt/project-files/literature/literature_table.xlsx (Papers and Gaps tabs).
% Limitations marked "Observed" in that table are the team's reading, not the authors'
% claims. They are written below as descriptions of what each study did, and should be
% checked against the full texts before submission.

\section{Introduction}
This chapter reviews the work this thesis builds on. Section~\ref{sec:lr-theory} sets out
the volatility models and the notion of an anomaly used throughout. Section~\ref{sec:lr-themes}
reviews seven bodies of work: hybrid GARCH and deep learning models, outliers in GARCH
estimation, volatility studies of the Dhaka Stock Exchange (DSE), anomaly detection in
financial markets, deep time-series anomaly detection, the evaluation of such detectors,
and transfer learning in finance. Section~\ref{sec:lr-gaps} lists the gaps that motivate
this study.

\section{Theoretical Background}\label{sec:lr-theory}

\subsection{Volatility Clustering and GARCH-Family Models}
Daily stock returns are close to uncorrelated, but their magnitude is not: large moves
tend to follow large moves. Engle's ARCH model~\cite{engle1982autoregressive} captured this
by letting the conditional variance depend on past squared innovations, and
Bollerslev's GARCH model~\cite{bollerslev1986generalized} added a dependence on past
conditional variance, which gives a parsimonious description of persistent volatility.
Two extensions allow negative shocks to raise volatility more than positive shocks of the
same size: EGARCH~\cite{nelson1991conditional}, which models the logarithm of the variance,
and the GJR model~\cite{glosten1993relation}, which adds a term that is active only after
negative returns. Chapter~\ref{ch:method} gives the equations used in this thesis.

Dividing each return innovation by its conditional standard deviation gives a standardized
residual, $z_t = \varepsilon_t/\sigma_t$. If the volatility model is adequate, $z_t$ has no
remaining volatility clustering, so a day with a large $|z_t|$ is unusual relative to the
volatility of its own period rather than relative to the whole sample.

\subsection{Anomalies in Financial Time Series}
\todo{Define point, contextual and collective anomalies, and state which kind this thesis
targets. Golmohammadi and Zaiane~\cite{golmohammadi2015time} is a starting point for
contextual anomalies in stock prices.}

\section{Review of Key Themes}\label{sec:lr-themes}

\subsection{GARCH Models Combined with Deep Learning}
Kim and Won~\cite{kim2018forecasting} fed the outputs of GARCH, EGARCH and EWMA models
into an LSTM and forecast the realized volatility of the KOSPI~200 index, reporting lower
errors than either the econometric models or the networks alone. Koo et
al.~\cite{koo2022hybrid} combined GARCH-type inputs with a root-type transform of the
volatility distribution before an LSTM, motivated by the observation that volatility is
concentrated near zero, so a network trained on the raw distribution fits the common
region and misses the tail. They report a 21\% reduction in RMSE over earlier LSTM-GARCH
hybrids on S\&P~500 realized volatility.

Both studies use GARCH output to improve a volatility forecast on a single index. Neither
addresses anomaly detection. The second study's motivation is relevant here: if a network
trained on raw values under-fits the tail, a transform that removes the volatility level
may help a detector separate unusual days from merely turbulent ones. This thesis tests
that idea with GARCH standardization.

\subsection{Outliers in GARCH Estimation}
Charles and Darn{\'e}~\cite{charles2005outliers} showed that outliers in daily returns
bias GARCH parameter estimates and test statistics, and proposed a procedure to detect and
correct them during estimation. Their work treats outliers as a nuisance for estimation
rather than as targets. For this thesis it has two consequences. Labelling a day as
anomalous from an extreme standardized residual is an established idea, and extreme days in
the estimation window will distort the fitted model, which is one reason the volatility
models here are estimated on the training block only.

\subsection{Volatility Studies of the Dhaka Stock Exchange}
Siddikee et al.~\cite{siddikee2016volatility} fitted GARCH(1,1) and ARCH models to the DSE
General Index (DGEN) from January 2002 to July 2013 and identified an episode of abnormal
volatility between 2009 and 2012, peaking in the first half of 2011. They report that
GARCH could not be applied in some years and used ARCH instead. Haque et
al.~\cite{haque2021modeling} fitted long-memory models (FIGARCH, FIEGARCH, FIAPARCH and
HYGARCH) to the DSEX and DSES indices and examined asymmetry through news impact curves.

Both studies are in-sample descriptions of DSE volatility. Neither evaluates models out of
sample or addresses anomaly detection. The DGEN series studied by Siddikee et al.\ was
replaced by DSEX in 2013, so their sample and the DSEX series used in this thesis do not
overlap.

\subsection{Anomaly Detection in Financial Markets}
Two surveys frame the area. Ahmed et al.~\cite{ahmed2016survey} review anomaly detection
for financial fraud, with an emphasis on clustering methods, and Hilal et
al.~\cite{hilal2022financial} review supervised, unsupervised and semi-supervised fraud
detection, including deep learning, and identify the scarcity of public labelled data as a
recurring obstacle.

Work on stock markets faces the same labelling problem. Golmohammadi and
Zaiane~\cite{golmohammadi2015time} detected contextual anomalies in S\&P~500 constituents
by comparing each stock with its peer group, and evaluated on injected synthetic anomalies
because confirmed manipulation cases were not available. On the Stock Exchange of Thailand,
Leangarun et al.~\cite{leangarun2021stock} trained an autoencoder and a GAN on normal
limit-order-book behaviour and detected five of six prosecuted manipulation cases, and
Chullamonthon et al.~\cite{chullamonthon2023ensemble} stacked an LSTM classifier and an
LSTM autoencoder with an SVM and detected all six. Six cases are too few to support
statistical comparisons between methods. Mariani et al.~\cite{mariani2021autoencoder}
applied LSTM, bidirectional LSTM and convolutional autoencoders to the IBOVESPA index and
flagged days with high reconstruction error, without a defined ground truth.

None of these studies adjusts its inputs for volatility clustering, and each is confined to
a single market.

\subsection{Deep Time-Series Anomaly Detection Methods}
Most deep detectors for time series score a point by how poorly a model trained on normal
data reconstructs or predicts it. Malhotra et al.~\cite{malhotra2016lstm} used an LSTM
encoder-decoder and scored windows by reconstruction error on sensor and medical data.
Hundman et al.~\cite{hundman2018detecting} forecast spacecraft telemetry with an LSTM and set
thresholds on the forecast errors with a nonparametric dynamic method. Xu et
al.~\cite{xu2018unsupervised} trained a variational autoencoder (Donut) on sliding windows
of web-service KPIs. Xu et al.~\cite{xu2022anomaly} proposed the Anomaly Transformer, which
scores points by an association-discrepancy criterion.

These methods were developed and tested on sensor, telemetry and operations data. None was
evaluated on financial returns, whose variance changes over time. A reconstruction-error
detector trained on raw returns will tend to score every day of a high-volatility period as
anomalous, which is the behaviour this thesis aims to measure and reduce.

\subsection{Evaluation of Time-Series Anomaly Detection}
Several papers question how these detectors are evaluated. Xu et
al.~\cite{xu2018unsupervised} introduced point adjustment: if any point inside a labelled
anomalous segment is detected, every point in the segment is counted as detected. Kim et
al.~\cite{kim2022towards} showed that under point adjustment, random anomaly scores can
match published state-of-the-art results, and proposed PA\%K, which credits a segment only
if at least $K$\% of its points are detected. Tatbul et al.~\cite{tatbul2018precision}
defined range-based precision and recall with tunable terms for existence, overlap,
position and cardinality. Paparrizos et al.~\cite{paparrizos2022volume} proposed the
threshold-free VUS-ROC and VUS-PR measures, which tolerate small lags between predicted and
true anomalies.

The benchmarks themselves have also been criticised. Wu and
Keogh~\cite{wu2023current} argue that widely used benchmarks are often trivial, have
unrealistic anomaly densities and contain labelling errors, and they released the UCR
anomaly archive. TSB-UAD~\cite{paparrizos2022tsb} and the evaluation by Schmidl et
al.~\cite{schmidl2022anomaly} compare many detectors on large collections of real and
synthetic series. Schmidl et al.\ report that no single algorithm wins and that simple
methods are often competitive. Neither collection includes financial series.

This thesis therefore reports event-based F1 alongside point-adjusted F1 and PA\%K, and
includes simple statistical baselines.

\subsection{Transfer Learning in Finance}
Koshiyama et al.~\cite{koshiyama2022quantnet} trained QuantNet, which shares a
market-agnostic layer across equity markets and adds market-specific encoders and decoders,
and measured gains in trading performance (Sharpe and Calmar ratios). Nguyen and
Yoon~\cite{nguyen2019novel} pretrained an LSTM on a large pool of stocks and fine-tuned it to
predict short-term price movements of target stocks. Both target prediction or trading
rather than detection, and neither transfers into a frontier market such as Bangladesh.

\section{Critical Analysis and Research Gaps}\label{sec:lr-gaps}
The review points to seven gaps.

\begin{enumerate}
    \item Hybrid GARCH and deep learning models are used for volatility forecasting
    \cite{kim2018forecasting,koo2022hybrid}. None of the studies reviewed uses
    GARCH-standardized returns as the input to a deep anomaly detector.
    \item Anomaly detectors applied to stock data do not adjust for volatility
    clustering~\cite{golmohammadi2015time,mariani2021autoencoder}, so calm and turbulent
    periods are scored on the same scale.
    \item Published DSE studies fit GARCH-family models in
    sample~\cite{siddikee2016volatility,haque2021modeling}. None of the studies reviewed
    applies deep or out-of-sample anomaly detection to the DSE.
    \item Cross-market transfer has been studied for trading and price
    prediction~\cite{koshiyama2022quantnet,nguyen2019novel}, not for anomaly detection into a
    frontier market.
    \item Labels in financial anomaly studies are either very few real cases
    \cite{leangarun2021stock,chullamonthon2023ensemble} or synthetic
    \cite{golmohammadi2015time}, and the shortage of public labelled data is a known
    obstacle~\cite{hilal2022financial}.
    \item The inflation caused by point adjustment has been measured on industrial and
    operations benchmarks~\cite{kim2022towards,wu2023current} but not on financial data.
    \item Large anomaly detection benchmarks contain no financial
    series~\cite{paparrizos2022tsb,schmidl2022anomaly}.
\end{enumerate}

This thesis addresses the first five directly. It compares an autoencoder on raw and on
GARCH-standardized returns, analyses false alarms by volatility regime, evaluates on DSE
data with time-based splits, transfers detectors from global indices to the DSE, and builds
labels from three sources: known market events, extreme standardized residuals and injected
anomalies. It addresses the sixth by reporting point-adjusted and unadjusted scores side by
side.

\section{Summary of Literature Review}
GARCH-family models describe DSE volatility well in sample, and deep detectors work on
sensor and operations data, but the two have not been combined for anomaly detection in
stock markets, and nothing comparable exists for the DSE. Evaluation practice in the field
also overstates performance when point adjustment is used. Chapter~\ref{ch:method} describes
the data, labels, models and evaluation protocol designed around these gaps.
